\section{From script to report}

\begin{frame}{Why use a rendered report?}
A script can compute a result. A policy brief must also tell the reader:
\begin{itemize}
\item the question and decision context;
\item the source, sample, unit, and period;
\item how the table and figure were produced;
\item what the evidence supports; and
\item what remains uncertain.
\end{itemize}
\vspace{0.3cm}
Quarto keeps the prose, code, and generated output in one document.
\end{frame}

\begin{frame}{R Markdown and Quarto}
\small
\begin{tabularx}{\textwidth}{p{2.6cm} X X}
\toprule
 & R Markdown & Quarto \\
\midrule
Source file & \texttt{.Rmd} & \texttt{.qmd} \\
Code chunks & R and other engines & R and other engines \\
Typical outputs & HTML, PDF, Word & HTML, PDF, Word, slides, books \\
Course use & Existing reports remain valid & New Math Camp briefs use Quarto \\
\bottomrule
\end{tabularx}
\vspace{0.35cm}
Both combine narrative and code. We use Quarto because one syntax supports the student book, slides, and standalone briefs.
\end{frame}

\begin{frame}[fragile]{Anatomy of a Quarto file}
\begin{lstlisting}
---
title: "Income and under-five mortality"
format:
  html:
    toc: true
    embed-resources: true
---

## Data and sample

```{r}
library(tidyverse)
health <- read_csv("data/lesson-4/health-2022.csv")
```
\end{lstlisting}
The YAML controls the document. Markdown supplies prose and headings. R chunks compute the evidence.
\end{frame}

\begin{frame}{Live demonstration: build the health brief}
Open \texttt{materials/lesson-4-demo.qmd} beside RStudio.
\begin{enumerate}
\item Read the question and sample statement.
\item Run the import and inspection chunk.
\item Run the grouped table.
\item Run the figure.
\item Compare the prose with the output.
\item Click \textbf{Render}.
\end{enumerate}
\vspace{0.25cm}
We will deliberately change one object name and render again to observe how a broken dependency appears.
\end{frame}

\begin{frame}{Rendering checks the document from the beginning}
When you click Render, Quarto starts a fresh R session and executes the chunks in order.
\begin{itemize}
\item An object created only in the Console is unavailable.
\item A missing package or file path becomes visible.
\item A table or figure must be generated by code in the document.
\item A successful render produces a portable record of the analysis.
\end{itemize}
\end{frame}

\begin{frame}{Checkpoint 2: repair the evidence chain}
A draft brief contains:
\begin{itemize}
\item a table restricted to 2022;
\item a figure pooling all years from 2000 to 2022; and
\item the sentence, ``Higher income reduces under-five mortality.''
\end{itemize}
\vspace{0.25cm}
Identify two problems and propose the smallest defensible revision to the table, figure, or sentence.
\vfill
\centering\textbf{2 minutes discuss \(\longrightarrow\) 2 minutes share}
\end{frame}

\begin{frame}{Checkpoint 2: resolution}
\begin{enumerate}
\item The table and figure use different periods and samples. A reader cannot treat them as evidence for the same comparison without explanation.
\item The sentence is causal. The cross-sectional comparison supports an association, not the effect of an income intervention.
\end{enumerate}
\vspace{0.3cm}
Smallest revision: align both outputs to 2022 and describe the observed association and its variation.
\end{frame}

\begin{frame}{The agent enters after the first interpretation}
Save the table, figure, finding, and limitation before opening Codex.
\vspace{0.3cm}
\textbf{The read-only request asks Codex to:}
\begin{enumerate}
\item reproduce the table from the stated sample;
\item inspect the variables, units, and figure mappings;
\item challenge one consequential statement with exact evidence; and
\item end with accept, revise, or reject and at most one revision.
\end{enumerate}
\vspace{0.2cm}
The complete prompt is in the Lab 4 handout. Independently verify one material statement in R.
\end{frame}
